\section{Limitations}

Simulation only. The hardware phase was pre-gated on this evaluation
supporting the mechanism, and the gate closed; every number here is
conditioned on the simulator's service model. The model's constants are
literature-anchored, marked low confidence, and their measured sensitivity
is the paper's principal caveat: the binding constant moves the
representative deficit between 4.4 and 21.5 percent, sign-stable
throughout. Synthetic workloads: no public traces of long-horizon agentic
serving existed at specification freeze; generators, anchors, and seeds
are released, and the adversarial family is constructed, not observed.
Three seeds per cell (Amendment 3): intervals are wider at low load, and
several 0.7x cells have upper bounds crossing zero individually, though
the verdict is invariant to any reading of them, since zero cells reach
the support threshold under any interpretation. Cluster sizes 8 to 32:
the 64-device axis was removed for measured infeasibility at the
corrected horizon; the trend across measured sizes is flat, and the
memory-bound regime is probed at one configuration rather than swept.
Value weights are operator-supplied, fixed at registration, and swept in
sensitivity. H2 and H3 are measured on representative configurations, not
crossed with the full grid. The two-cell ablation design cannot separate
any element's effect from its measured seed-noise floor at three
replications; the campaign attributes the mechanism's cost to the
workload's structural duty cycle, and finer per-element attribution
requires swept axes and more seeds (data released). Finally, the
registration's one-fix clause means any residual measurement artifact, in
either direction, stands uncorrected in these numbers by design; the
amendment log is the complete record of what was corrected and when.
